% Editable transcription; verify reused text against the original PDF.
\section*{Introduction}

In 1980 Roger C. Schank wrote an editorial for Intelligence entitled
``How much intelligence is there in artificial intelligence (AI)?''.
His first observation was the lack of any interaction between the fields
of intelligence research and artificial intelligence research.
Since then limited interactions have taken place.
His second observation was that AI is relevant for intelligence research.
This was based on the state of the art of research in both fields, at that
time, when AI was still in its infancy.
Given the breathtaking developments in modern AI, it is worthwhile to ask
Schank's question again.
We contend that the question's relevance has increased over time.
Shank's third observation was that real intelligence is all about
generalization, which at that time was a weak point in AI, but perhaps
not anymore.

Schank was modest in claiming to speak for AI, given the multitude and
diversity of AI approaches at that time.
The current explosion of developments in AI research makes a systematic
overview even more impossible.
But we think it is useful to present a short overview here and discuss
the relations between AI and intelligence research as it is reported
in Intelligence.

Perhaps it is helpful to start discussing possible reasons why AI may
be irrelevant to intelligence research.
The first well known argument is that AI may be building intelligent
machines, but that machine intelligence is not similar to human
intelligence or cognition.
Dennett's (2006) cognitive wheel idea illustrates this.
The wheel is a beautiful invention, but it is not used in nature as a
method of movement and therefore of no interest to biologists.
This may apply to many AI inventions, e.g., cognitive wheels.
The second reason is that AI, even if it has psychological relevance,
can only inform us about cognition and not intelligence, as the latter
primarily concerns individual differences in cognitive functioning,
and not cognitive functioning itself.
This relates to Cronbach's famous distinction between the psychology
of processes or mechanisms versus the psychology of individual
differences, where both apparently exist independently of each other
(Cronbach, 1957).
We will argue against both arguments, but first discuss current
developments in AI.
